% =====================================================================
%  Vorlage: Betriebsanweisung für einen Gefahrstoff
%  Kopieren nach betriebsanweisung/ba_<stoff>.tex und ausfüllen.
%  Grundlage: Sicherheitsdatenblatt (SDB) des Herstellers, das im FabLab
%  vorliegen muss. Nummer: BA-GS-<Nr.> (fortlaufend für das ganze FabLab).
%  GHS-Piktogramme: GHS01 ... GHS09 (fablab-document/ba-symbole/QUELLEN.md)
% =====================================================================
\BAsetup{
  typ=gefahrstoff,
  titel={Handelsname -- chemische Bezeichnung},
  kurztitel={Gefahrstoff: Handelsname},
  taetigkeit={Tätigkeiten mit dem Stoff},
  nummer=BA-GS-XX,
  stand=TT.MM.JJJJ,
  verantwortlich={Name der verantwortlichen Person},
  signalwort=GEFAHR,                  % oder ACHTUNG, laut SDB
  entwurf,                            % entfernen, sobald freigegeben
}

\BAkopf

\BAblock{GEFAHREN FÜR MENSCH UND UMWELT}{GHS07}{%
  \begin{itemize}
    \item Gefahrenhinweise (H-Sätze) laut Sicherheitsdatenblatt.
  \end{itemize}}

\BAblock{SCHUTZMASSNAHMEN UND VERHALTENSREGELN}{M004,M009,P022}{%
  \begin{itemize}
    \item Schutzausrüstung laut Sicherheitsdatenblatt.
    \item Essen und Trinken verboten. Nach der Arbeit Hände waschen.
  \end{itemize}}

\BAblock{VERHALTEN IM GEFAHRFALL}{W001,F001}{%
  \begin{itemize}
    \item Verschütteten Stoff aufnehmen, Betreuer informieren.
    \item Brand: geeignete Löschmittel laut SDB. \textbf{Feuerwehr: 112}
  \end{itemize}}

\BAblock{ERSTE HILFE}{E003,E011}{%
  \begin{itemize}
    \item \textbf{Augen:} ...
    \item \textbf{Haut:} ...
    \item \textbf{Einatmen:} ...
    \item \textbf{Verschlucken:} ...
  \end{itemize}\vspace{1.5mm}
  \textbf{Notruf: 112}\qquad Giftnotruf München: 089\,/\,19240\qquad
  FabLab: 09131\,/\,85\,28013}

\BAletzterblock{SACHGERECHTE ENTSORGUNG}{}{%
  \begin{itemize}
    \item Entsorgung laut Sicherheitsdatenblatt.
  \end{itemize}}
